\section{Finite orbit budgets and exact predecessor catalogues}
\label{sec:v7-finite-core}

This section records finite arithmetic statements in the post-v6 Lean
development. Their proofs do not import a predecessor-density theorem,
an atom floor, a mixing estimate, or a conjecture about divergent orbits.
The parity-cylinder geometry is classical
\cite{terras1976,lagarias2003}. Uniform reciprocal and correction-product
bounds for finite distinct trajectories are already proved in
\cite[Section~2]{valbuena2026}. A related separation by odd-step count
appears in Curry's collision-free window estimate
\cite[Theorem~2.3]{curry2026}. The contribution asserted here is the
specific parity-moment proof, its explicit rational constant, and its
integration into the local formal development. Qualitative Green
finiteness is not claimed as new.

Write
\[
 T(n)=\begin{cases}n/2,&n\text{ even},\\(3n+1)/2,&n\text{ odd},\end{cases}
 \qquad
 S(n)=\frac{3n+1}{2^{\nu_2(3n+1)}}\quad(n>0\text{ odd}).
\]
For an actual Syracuse segment $x_i=S^i(u)$, put
\[
 A_k=\sum_{i<k}\nu_2(3x_i+1),\qquad
 \Omega_k=\frac{3^k}{2^{A_k}}.
\]
For an exact exponent word $w$, the corresponding seed-independent gain
is $\Omega(w)=3^{|w|}/2^{\sum w}$. An exact word includes the prescribed
valuations and the odd terminal state, rather than only an affine identity.
A first arrival at $n$ at time $k$ means $x_k=n$ and $x_i\ne n$ for
every $i<k$. The odd-unit restriction $3\nmid u$ used in the subsequent
research programme is not needed for most estimates below; where only
positive odd seeds are specified, multiples of three are allowed as seeds.

\subsection{An explicit reciprocal budget for finite injective orbits}

\begin{theorem}[Finite reciprocal budget]\label{thm:v7-reciprocal}
Let $u>0$ and $L\ge0$. If $T^0(u),\ldots,T^{L-1}(u)$ are distinct, then
\begin{equation}\label{eq:v7-reciprocal}
 \sum_{i=0}^{L-1}\frac1{T^i(u)}
 \le\frac{133547}{1655}<81.
\end{equation}
No assumption is made about the orbit after time $L-1$.
\end{theorem}
\begin{proof}
Let $J_N(s)$ count the odd entries among
$T^0(s),\ldots,T^{N-1}(s)$. The exact parity lift and its moment identity
are
\begin{align}
 T^N(s+2^N)&=T^N(s)+3^{J_N(s)},
 &J_N(s+2^N)&=J_N(s),\label{eq:v7-parity-lift}\\
 \sum_{0\le s<2^N}q^{J_N(s)}&=(1+q)^N.\label{eq:v7-parity-moment}
\end{align}
The first identity follows by induction through the $N$ exposed parity
bits. The last newly exposed bit is complemented by adding $2^N$;
pairing the two lifts therefore multiplies the preceding moment by
$1+q$, proving the second identity. Here the auxiliary representative
$s=0$ is permitted, with $T(0)=0$.
Induction also gives
\[
 2^N T^N(s)<3^{J_N(s)}(s+2^N),
 \qquad 0\le s<2^N,
\]
so $T^N(s)<2\cdot3^{J_N(s)}$.

Consider a finite set $B$ in $[2^N,2^{N+1})$ on which $T^N$ is
injective. Write $n=2^N+s$. At a fixed value $J_N(s)=m$, distinct $n$
give distinct reduced endpoints $T^N(s)$ by
\eqref{eq:v7-parity-lift}. There are at most $2\cdot3^m$ such
endpoints. Consequently the number of elements with $J_N(s)<M$ is
less than $3^M$. On the other hand, substituting $q=3/2$ in
\eqref{eq:v7-parity-moment} bounds the number with $J_N(s)\ge M$ by
$(5/2)^N(2/3)^M$. Thus
\[
 \frac{\#B}{2^N}\le
 \frac{3^M}{2^N}+\left(\frac54\right)^N\left(\frac23\right)^M.
\]

Let $b_N$ count the states of the segment in the $N$th dyadic band.
Delete its last $N$ time indices, or the whole segment if it has fewer
than $N$ entries. The map $T^N$ is injective on the remaining states:
their images are still distinct states of the original segment.
Restoring the deleted indices costs at most $N$. With
$M=\lfloor3N/5\rfloor$, it follows that
\begin{align}
 \frac{b_N}{2^N}&\le\frac{N}{2^N}+h_N+\ell_N,
 \label{eq:v7-band-bound}\\
 h_N&=\left(\frac54\right)^N\left(\frac23\right)^{\lfloor3N/5\rfloor},
 \qquad \ell_N=\frac{3^{\lfloor3N/5\rfloor}}{2^N}.\notag
\end{align}
Only finitely many bands are occupied. Their reciprocal contribution is
at most $\sum_N b_N/2^N$. Grouping the upper envelopes into blocks of
five gives
\[
 h_{5r+s}=\left(\frac{3125}{3456}\right)^r h_s,
 \qquad
 \ell_{5r+s}=\left(\frac{27}{32}\right)^r\ell_s
 \quad(0\le s<5),
\]
with $\sum_{s<5}h_s=3271/576$ and
$\sum_{s<5}\ell_s=51/16$. Finite geometric sums therefore give the
respective bounds $19626/331$ and $102/5$. The terminal error costs at
most $\sum_N N/2^N\le2$. At $N=0$, however, the actual band contains
only the possible state $1$, whereas \eqref{eq:v7-band-bound} assigns
two units. Saving that extra unit yields
\[
 \sum_{i<L}\frac1{T^i(u)}
 \le 2+\frac{19626}{331}+\frac{102}{5}-1
 =\frac{133547}{1655}.
\]
All estimates used in the implementation are finite rational inequalities;
the displayed geometric bounds do not require an asymptotic counting
hypothesis.
\end{proof}

The principal declaration is
\ttbr{CollatzShadowing.OrbitReciprocal.finite\_orbit\_reciprocal\_sum\_le}.
The companion real-scalar theorem is
\ttbr{finite\_orbit\_reciprocal\_sum\_real\_lt\_eighty\_one}
in the same namespace. The temporal deletion, which is essential for the
injectivity assertion used above, is formalized in
\texttt{FiniteOrbitTransfer.lean}; injectivity of the full Collatz map on
arbitrary integer sets is not assumed.

\subsection{First arrivals and the additive correction product}

\begin{proposition}[First-arrival envelope]\label{prop:v7-first-arrival}
For every positive odd first-arrival segment $u\to n$ of length $k$,
\begin{equation}\label{eq:v7-product}
 n=\Omega_k u \mathcal P_k,
 \qquad
 \mathcal P_k=\prod_{i<k}\left(1+\frac1{3x_i}\right),
 \qquad 1\le \mathcal P_k\le P_0:=\exp(27).
\end{equation}
In particular, $\Omega_k u\le n\le P_0\Omega_k u$.
\end{proposition}
\begin{proof}
A finite trajectory through its first arrival at a target is injective,
including its endpoint. Indeed, a repetition before the endpoint would
make the remaining future repeat and would produce an earlier visit to
that same target. Every Syracuse step at an odd state has positive
exponent, and the elapsed shortcut time at $S^i(u)$ is exactly $A_i$.
All strictly interior shortcut states of such an episode are even.
Hence a first Syracuse arrival at the odd target $n$ is a first shortcut
arrival at time $A_k$. Its shortcut segment is injective, and the
Syracuse states $x_i$, $i<k$, are a subset of that segment. Theorem
\ref{thm:v7-reciprocal} gives $\sum_{i<k}1/x_i<81$.

Multiplying the exact identities
$x_{i+1}=(3/2^{a_i})x_i(1+1/(3x_i))$ proves
\eqref{eq:v7-product}. Each factor is at least one, and
$1+t\le e^t$ gives
$\mathcal P_k\le\exp(\sum_{i<k}1/(3x_i))\le e^{27}$.
\end{proof}

The first-arrival sum and exact shortcut-clock interface are the
declarations
\ttbr{FirstArrivalBudget.first\_arrival\_reciprocal\_sum\_lt\_eighty\_one}
and \ttbr{FirstArrivalBudget.shortcut\_clock}, under the common
\texttt{CollatzShadowing} prefix. The product identity and its exponential
estimate are
\ttbr{WordProductEnvelope.inflation\_eq\_product} and
\ttbr{WordProductEnvelope.inflation\_le\_exp\_of\_reciprocal\_budget}.
The identity itself allows repeated visits; the uniform constant $P_0$
here is justified by the first-arrival hypothesis. It must not be
reused for arbitrarily many turns around a cycle.

\begin{corollary}[Height of an existing recoding]
\label{cor:v7-height}
Suppose $w:u\to n$ is as in Proposition~\ref{prop:v7-first-arrival}
and $w':v\to y$ is another exact positive odd-source word. If
$\Omega(w')\ge\Omega(w)/C$ for $C>0$, then
\begin{equation}\label{eq:v7-height}
 v\le C\frac{y}{n}P_0u.
\end{equation}
In particular, equal word lengths and equal exponent totals allow $C=1$.
If additionally $y\le21n$, then $v\le21e^{27}u$.
\end{corollary}
\begin{proof}
The nonnegative affine constant of the exact word $w'$ gives
$\Omega(w')v\le y$. The first-arrival product for $w$ gives
$1/\Omega(w)=u\mathcal P_k/n\le P_0u/n$. Combining the two inequalities proves
\eqref{eq:v7-height}.
\end{proof}

This is implemented by
\ttbr{WordProductEnvelope.comparable\_slope\_height} and, for the
same-profile specialization,
\ttbr{FirstArrivalBudget.first\_arrival\_same\_profile\_height}.
The canonical unit sibling of a unit $n$ is $4n+1$ if it is a unit,
and $16n+5$ otherwise; both possibilities are at most $21n$.
The canonical-sibling specialization is
\ttbr{FirstArrivalBudget.canonical\_first\_arrival\_profile\_height}.
These are bounds on a recoding \emph{that exists}. They supply neither
that recoding nor a positive proportion of successfully recoded sources.

\subsection{Squared occupation, including repeated cyclic visits}

Let $\phi_k=\Omega_k/x_k$. The exact decrement of its square is
\begin{equation}\label{eq:v7-square-loss}
 \phi_k^2-\phi_{k+1}^2
 =\Omega_k^2\frac{6x_k+1}{x_k^2(3x_k+1)^2}.
\end{equation}
This follows directly from
$\phi_{k+1}=3\Omega_k/(3x_k+1)$ and does not require first arrival.
Put
\[
 V(n)=\frac{n^2(3n+1)^2}{6n+1}\quad(n>0).
\]

\begin{theorem}[Finite quadratic occupation]\label{thm:v7-square}
For a fixed positive odd seed $u$, a target $n>0$, and any finite set
$H$ of times with $S^k(u)=n$ for all $k\in H$,
\begin{equation}\label{eq:v7-visit-bound}
 \sum_{k\in H}\Omega_k^2\le\frac{V(n)}{u^2}.
\end{equation}
For any finite catalogue $\mathcal W$ of distinct exact words ending at
$n$, with arbitrary positive odd sources, even when sources repeat,
\begin{equation}\label{eq:v7-catalogue-square}
 \sum_{w\in\mathcal W}\Omega(w)^2\le2V(n).
\end{equation}
Periodic targets and the empty word are allowed.
\end{theorem}
\begin{proof}
The nonnegative losses \eqref{eq:v7-square-loss} telescope from
$\phi_0^2=1/u^2$. At a visit to $n$, the loss is exactly
$\Omega_k^2/V(n)$, proving \eqref{eq:v7-visit-bound} for every finite
set of visiting times. From one fixed source there is only one actual
word of any specified length; distinct words in its fibre therefore
correspond to distinct visiting times, including returns. Group the
catalogue by source and use
\[
 \sum_{u\in U}\frac1{u^2}\le2
\]
for any finite set $U$ of positive integers. The latter estimate follows,
for example, by telescoping $1/u^2\le2/u-2/(u+1)$.
\end{proof}

The exact loss, repeated-visit bound, and arbitrary-source catalogue
bound are respectively
\ttbr{QuadraticOccupation.potential\_sq\_sub\_succ},
\ttbr{QuadraticOccupation.sum\_weight\_sq\_visits\_le}, and
\ttbr{QuadraticOccupation.tilted\_catalogue\_le}.
The simpler fixed-depth estimate $\sum\Omega(w)^2\le2n^2$ is
\ttbr{IntegerAtomTilt.tilted\_integer\_atom\_le}; it uses source
injectivity at that depth. The same constant is shared across arbitrary
depths if the target has no positive return, as stated by
\ttbr{IntegerAtomTilt.tilted\_nonreturn\_catalogue\_le}.
At a cyclic target this additional hypothesis cannot be dropped from
the $2n^2$ argument; Theorem~\ref{thm:v7-square} supplies the separate
repeated-visit estimate.

The corresponding infinite integer Green sum was already shown finite
in \cite[Proposition~11]{weaver9v1}. Equation
\eqref{eq:v7-catalogue-square} is a finite elementary realization with
an explicit constant; it is not a priority claim for Green finiteness.

\begin{corollary}[Finite bounded-slope catalogue]
\label{cor:v7-bounded-catalogue}
For every positive integer target $n$ and positive rational $R$, the
set of all exact words from positive odd sources to $n$ satisfying
$\Omega(w)\ge1/R$ is finite, and
\begin{equation}\label{eq:v7-catalogue-card}
 \#\mathcal W_R(n)\le2V(n)R^2.
\end{equation}
Every corresponding source satisfies $u\le Rn$. There exist integers
$K,A$ bounding all their word lengths and exponent totals, respectively.
The conclusion includes cyclic targets.
\end{corollary}
\begin{proof}
Each selected word contributes at least $1/R^2$ to
\eqref{eq:v7-catalogue-square}. Hence every finite subcatalogue has
cardinality at most the fixed number $2V(n)R^2$; an infinite set would
have finite subsets of arbitrarily large cardinality, a contradiction.
The source bound follows from $\Omega(w)u\le n$. The length and total
bounds are maxima of two functions on the resulting finite set.
\end{proof}

The five public results in \texttt{BoundedSlopeCatalogue.lean} are
\ttbr{subcatalogue\_card\_le}, \ttbr{catalogue\_finite},
\ttbr{catalogue\_ncard\_le}, \ttbr{catalogue\_parameter\_bounds}, and
\ttbr{catalogue\_source\_bound}. Its implemented catalogue is slightly
larger than the odd-source catalogue displayed here: only positivity of
the source is imposed. The numerical source and cardinality bounds are
explicit. The proof supplies the temporal bounds $K,A$ only
\emph{existentially}; neither a numerical hitting-time cutoff nor a
procedure whose stopping criterion follows from the cardinality bound
is asserted.

\subsection{Finite horizons and the direction of collision certificates}

Two additional interfaces delimit what these finite upper bounds can
and cannot provide. First, every positive odd segment obeys
\begin{equation}\label{eq:v7-shifted-growth}
 2^k(S^k(u)+1)\le3^k(u+1).
\end{equation}
Thus a word of length at most $K$ ending at $n$ has source at least
$R\ge1$ whenever $3^K(R+1)\le2^K(n+1)$. For a finite catalogue with
distinct sources at least $R$,
\[
 \sum_w\Omega(w)^2\le n^2\sum_u\frac1{u^2}
 \le\frac{2n^2}{R}.
\]
Distinctness is automatic at one fixed depth, and across arbitrary
depths at a target without positive return. These are
\ttbr{FiniteHorizonBudget.iterate\_shifted\_growth},
\ttbr{tilted\_tail\_catalogue}, \ttbr{tilted\_fixed\_depth\_tail},
and \ttbr{tilted\_nonreturn\_horizon}. Repeated sources at a cyclic
target require their own occupation budget; they are not silently
treated as distinct.

Second, let $R_0,\ldots,R_N$ be nondecreasing real numbers with
$R_i\ge a>0$. Suppose a finite subprobability catalogue $P_i$ contains
a distinguished branch with mass $R_i/R_{i+1}$, and let
$D_i=1-\sum_{j\in P_i}p_{i,j}^2$. Then
\begin{equation}\label{eq:v7-collision-budget}
 0\le D_i\le1-\left(\frac{R_i}{R_{i+1}}\right)^2,
 \qquad
 \sum_{i<N}D_i\le\frac2a(R_N-R_0).
\end{equation}
Indeed $1-r^2\le2(1-r)$, and
$(R_{i+1}-R_i)/R_{i+1}\le(R_{i+1}-R_i)/a$; the increments telescope.
The finite algebra is
\ttbr{SpineDispersion.collision\_spine\_budget}.

If two distinct actual branches each have mass at least $c$, every
branch has mass at most $1-c$, whence $D_i\ge c$. This is
\ttbr{SpineDispersion.collision\_defect\_ge\_of\_two\_branches};
the more general lower-allocation certificate is
\ttbr{collision\_defect\_ge\_lower\_allocations}. These statements
do not prove that such branches exist. Moreover, a truncated catalogue
has \emph{larger} defect than an extension, because adding squared
masses decreases $1-\sum p^2$. Its computed defect is therefore an
\emph{upper}, not a lower, bound on the full law's defect. A branch
lower certificate transfers only when it remains valid with the actual
unnormalized masses for every finite extension and the corresponding
limit is justified. Renormalized sample frequencies do not supply this
certificate.

The named declarations are tied to the repository's pinned verification
record. The seven-module integration on 15 November audited all 136
compiled declarations in that scope; the subsequent four-module
integration audited all 68 declarations, including five encoded-private
helpers, and passed an incremental full-library build of 3404 jobs.
The permitted transitive axioms in those audits were
\texttt{propext}, \texttt{Classical.choice}, and \texttt{Quot.sound}.
A separate Ubuntu replay reproduced the latter audit and finite
arithmetic outputs. Those recorded checks use Lean and Mathlib
\texttt{v4.29.1}; they are not a cold rebuild of all dependencies or a
replay with an independently implemented kernel. The release manifest
identifies the sources and evidence included in this version.
